% Generated draft: requires mathematical and source review.
\documentclass[11pt]{article}
\usepackage[margin=25mm]{geometry}
\usepackage{fontspec}
\setmainfont{Latin Modern Roman}
\usepackage{amsmath,amssymb,mathtools,unicode-math}
\setmathfont{Latin Modern Math}
\usepackage{enumitem,needspace,xurl,hyperref}
\setlength{\parindent}{0pt}
\setlength{\parskip}{5pt}
\setlength{\emergencystretch}{3em}
\newcommand{\N}{\mathbb{N}}
\newcommand{\Z}{\mathbb{Z}}
\newcommand{\Q}{\mathbb{Q}}
\newcommand{\R}{\mathbb{R}}
\newcommand{\C}{\mathbb{C}}
\newcommand{\F}{\mathbb{F}}
\newcommand{\A}{\mathbb{A}}
\newcommand{\PP}{\mathbb{P}}
\newcommand{\E}{\mathbb{E}}
\newcommand{\eps}{\varepsilon}
\newcommand{\dd}{\,\mathrm{d}}
\begin{document}
\section*{UnsolvedMath ID: 801122}
\textbf{Source problem code:} catalogue entry
\subsection{Definitions and mathematical statement}
Let \(N\) be a \(\mathrm{II}_1\) factor with normalized trace \(\tau\). For a nonprincipal ultrafilter \(\mathcal U\) on \(\N\), write
\[
N^{\mathcal U}=\ell^\infty(\N,N)\big/\{(x_j):\lim_{j\to\mathcal U}\|x_j\|_2=0\},
\qquad \|x\|_2=\tau(x^*x)^{1/2}.
\]
Identify \(N\) with its diagonal image in \(N^{\mathcal U}\), and equip the ultrapower with its induced trace. In either algebra, a Haar unitary is a unitary \(u\) satisfying \(\tau(u^m)=0\) for every \(m\in\Z\setminus\{0\}\), with the induced trace used in the ultrapower. Denote the Haar unitaries of an algebra \(A\) by \(\mathcal H(A)\), and put \([a,b]=ab-ba\).

For \(u,v\in\mathcal H(N)\), define \(u\sim_{N^{\mathcal U}}v\) if there are finitely many \(w_1,\ldots,w_r\in\mathcal H(N^{\mathcal U})\) such that
\[
[u,w_1]=[w_1,w_2]=\cdots=[w_{r-1},w_r]=[w_r,v]=0.
\]
Having a unique sequential commutation orbit means
\[
\forall u,v\in\mathcal H(N),\qquad u\sim_{N^{\mathcal U}}v.
\]
Thus the endpoints are required to belong to \(N\), while intermediate unitaries may belong to its ultrapower.

For an explicit edge-count convention, define
\[
p_{N,\mathcal U}(u,v)=\min\left\{k\in\Z_{\geq0}:\begin{array}{l}
\text{there exist }x_0,\ldots,x_k\in\mathcal H(N^{\mathcal U}),\\
x_0=u,\ x_k=v,\quad [x_{i-1},x_i]=0\ (1\leq i\leq k)
\end{array}\right\},
\]
with value \(+\infty\) if no such path exists. When there is a unique orbit, set
\[
\mathfrak O_{\mathcal U}(N)=\sup_{u,v\in\mathcal H(N)}p_{N,\mathcal U}(u,v).
\]
The source writes this diameter as \(\mathfrak O(M)\) although its defining supremum uses \(N\); the notation above uses \(N\) consistently.

Using the finite-set convention, property Gamma means that for every finite \(F\subseteq N\) and every \(\varepsilon>0\), there is a unitary \(z\in N\) such that
\[
\tau(z)=0,\qquad \|[z,a]\|_2<\varepsilon\quad(a\in F).
\]
A non-Gamma factor fails this property. The source does not impose separability or specify a convention for property Gamma in the nonseparable setting.

\textbf{Existence question, with unresolved parameters made explicit.}
For a prescribed finite \(n\in\N\), consider the assertion that there exists a \(\mathrm{II}_1\) factor \(N\) such that
\[
\begin{gathered}
N\text{ is non-Gamma},\\
\forall u,v\in\mathcal H(N),\quad u\sim_{N^{\mathcal U}}v,\\
\mathfrak O_{\mathcal U}(N)=n.
\end{gathered}
\]
The supplied question asks for a non-Gamma factor with a unique sequential commutation orbit and sequential diameter equal to \(n\in\N\). It does not explicitly specify whether existence is requested for every prescribed \(n\) or merely for some finite \(n\), whether \(0\in\N\), or how \(\mathcal U\) is quantified. The displayed assertion is therefore a parameterized formulation, not a resolution of those ambiguities. Moreover, the edge-count convention above requires confirmation against the cited source: counting intermediate unitaries instead would change the numerical diameter.
\subsection{Short English statement}
The question asks whether a non-Gamma \(\mathrm{II}_1\) factor can have all its Haar unitaries connected by finite chains of commuting Haar unitaries in an ultrapower, with the supremum of the shortest chain lengths exactly a specified finite \(n\). The quantifier over \(n\), the ultrafilter convention, the path-length normalization, and any separability hypothesis remain unresolved in the supplied record.
\subsection{Sources}
\begin{itemize}
\item Local UnsolvedMath registry: permanent ID \(801122\), exact source problem code \texttt{catalogue entry}, title ``Ultraproduct \(\mathrm{II}_1\) Factors --- Problem 23.'' The supplied record attributes the problem to Elayavalli, with proposed year 2025, and labels its status as open. That status is supplied registry metadata, not an independently verified mathematical or literature assessment. Upstream revision: \texttt{d371ee917ac9647fbb06cd6fc529deec6f3dda82}.
\item Elayavalli, \emph{50 Open Problems: Ultraproduct \(\mathrm{II}_1\) Factors} (2025), Problem 23. Supplied source URL: \url{https://arxiv.org/abs/2511.20377}. The provenance records access on 2026-10-16 and extraction by \texttt{arxiv\_latex}; these are supplied provenance claims, not a claim of independent source access or verification.
\item The supplied statement cites the reference label \texttt{elayavalli2023sequential} for sequential commutation. No full bibliographic entry or URL for that reference is supplied, so its definitions and path-length conventions have not been independently checked.
\end{itemize}
\end{document}
